\section{Dopants, color centers and small targets}\label{app:E.21}

A final state that the band picture does not contain occurs in donor clusters near the metal--insulator transition, molecular excitons with charge-transfer admixture, the ejected electron of a graphene sheet, the self-trapped exciton that becomes a Frenkel pair, and a molecule that dissociates in its cage (\cref{chap:16}). \Cref{tab:E21-params} summarizes the calculations.

\subsection{Disorder-averaged donor responses}

Phosphorus donors in silicon at $\nd_D = 10^{18}$ and $1.8\times10^{18}\,\si{cm^{-3}}$ lie at $\nd_D/\nd_c \approx 0.3$--0.6, where the dilute response of \cref{sec:16.5} fails at order one near the edge: donor pairs and clusters broaden the lines into an impurity band. The dilute response is replaced by the configurational average $\langle S^D(\vect{q},\omega)\rangle$ per donor, in \si{eV^{-1}}, the per-cell convention of \cref{app:E.17} with the donor in place of the cell, at $n_B = 0$ and $n_g = 0$ on $\mathfrak{T}$, with a Hartree mean field on the donor electrons and host screening. Each donor carries the six-valley orbitals of \cref{sec:16.5}, split into $A_1$, $E$ and $T_2$, at Poisson-random positions. The intervalley terms of the transition density lie at $\vect{q} = \vect{k}_\iota - \vect{k}_{\iota'} + \vect{G}$, between 0.685 and \SI{3.88}{keV}, and the lightest particle that delivers the \SI{11.70}{meV} of the $A_1 \to T_2$ line at \SI{0.685}{keV} has \SI{130}{keV}, so the $1s(A_1) \to 1s(E,T_2)$ lines, absent in optics, appear at finite $q$ for MeV-scale dark matter.

A Hubbard--Anderson tight-binding model of the random donor ensemble carries the disorder. Its hopping comes from the overlaps of the multivalley envelopes (effective Bohr radius $a^*$, with the central-cell correction on each site), and it is solved for $10^4$ to $10^5$ sites by Lanczos or the kernel polynomial method and averaged over 100 realizations, for about $10^3$ core-hours. DFT on 1000- to 4000-atom silicon clusters with two to four phosphorus atoms, at the shared \SI{60}{Ry} setting of \cref{app:E.17}, checks the central-cell part for about $10^4$ more. The ensemble has to satisfy the Thomas--Reiche--Kuhn sum per donor, reproduce the single-donor response of \cref{sec:16.5} in the dilute limit, and give $\langle S^D\rangle$ stable to \SI{5}{\percent} in the number of sites and realizations. The hopping oscillates with distance because of valley interference, so a single-valley envelope is wrong; compensation changes the occupation and must be stated. The Bloch functions at the conduction minima, on the same grid, also give the coefficients $\mathcal{U}_{j'j}(\vect{G})$ that set the strengths of the intervalley lines, and with the dielectric matrix of the host from \cref{app:E.6} they give the local-field ratio $\mathfrak{f}^{\mathrm{loc}}_{\mathrm{D}}$ of \cref{eq:16.5-floc}, which \cref{sec:16.5} takes as one.

\subsection{Molecular crystals}

In anthracene and trans-stilbene the lowest excitons are Frenkel-like, localized on one molecule, but they mix with charge-transfer states in which the electron is on a neighbor. A two-state model with transfer integral $t_{\mathrm{CT}}$ and energy splitting $E_{\mathrm{CT}} - E_{\mathrm{Fr}}$ gives a mixing angle $\tan2\theta_{\mathrm{CT}} = 2t_{\mathrm{CT}}/(E_{\mathrm{CT}} - E_{\mathrm{Fr}})$ (\cref{sec:16.4}), and the charge-transfer form factor carries the intermolecular phase, which no Frenkel model has; for orthogonal orbitals it vanishes at $q = 0$ and enters at order $q^2d^2$, with $d$ the intermolecular distance. The periodic Bethe--Salpeter response measures how much this matters: we tabulate its difference $\Delta S(\vect{q},\omega)$ from the zero-overlap Frenkel response, the product $\Omega S(\vect{q},\omega)$ on $|\vect{q}| \le \SI{30}{keV}$ times 590 directions, the charge-transfer weight $\theta_{\mathrm{CT}}^2$, and, as the summary, the largest relative change of $|f|^2$ over $q$ for each state. The periodic calculation is at $\mathcal{X}$ = BSE with intermolecular overlap, and the Frenkel model is its zero-overlap truncation.

It is $G_0W_0$ plus BSE (Yambo or BerkeleyGW) on the measured crystal structures from the Cambridge Structural Database, with 50 to 100 atoms per cell and a $4^3$ grid at finite $q$. Anthracene has 66 and trans-stilbene 136 occupied bands, so $G_0W_0$ needs about 1000 bands, or a Sternheimer formulation that needs no empty states; the BSE needs about 200. The molecular inputs of the Frenkel model come from TDDFT or molecular $GW$--BSE (PySCF or ORCA), and the Pariser--Parr--Pople model with standard parameters gives the semi-empirical comparison. The Frenkel code takes the transition densities on a real-space grid with $\Delta r \le \pi/q_{\mathrm{max}} \approx 0.4\,a_0$ at \SI{30}{keV}, Fourier transforms them, builds the intermolecular coupling by Ewald summation with the $\vect{G} = 0$ term removed (it belongs to $\varepsilon$), and diagonalizes the Frenkel Hamiltonian for about 30 states up to \SI{10}{eV}. The exciton energies must be stable to \SI{50}{meV} under doubling of bands and grid, the partial $f$-sum per molecule must hold to \SI{2}{\percent}, and $|f|^2$ at $q = \alpha m_e$ must be stable to \SI{5}{\percent} in the basis. Relaxed van der Waals geometries are not good enough; the BSE exchange term must carry local fields consistently; and the $\pi$-electron model fails at $q \gtrsim \SI{10}{keV}$, where the $\sigma$ electrons contribute. Each crystal costs $10^3$ to $10^4$ core-hours, and its tables about 10.

\subsection{Low-dimensional response tables}

Nanotubes, graphene sheets, layered stacks and quantum dots each screen with their own Coulomb kernel (\cref{sec:16.2}). For sheets and stacks we tabulate $S(q_\parallel,q_z,\omega)$, for tubes of radius $r_{\mathrm{t}}$ $S(q_\parallel,q_\perp,\omega)$, directional and per unit area of sheet or length of tube in the normalization of \cref{sec:4.2} (with $2\pi\delta$ in $S$), and for dots $S^{\mathrm{QD}}(\vect{q},\omega)$ in the dilute embedding of \cref{sec:16.3}. Two inputs are reported with them: the proper density response $\tilde R^{\mathrm{2D}}$ of the monolayer per unit area, from Yambo with a two-dimensional Coulomb cutoff or from a Wannier bubble, and the Keldysh length $r_0$ defined by its static remote-band part, $-\tilde R^{\mathrm{2D}}_{\mathrm{bg}} = 2r_0q_\parallel^2/e^2$ (\cref{sec:16.2}). The tables are at $\mathcal{X}$ = RPA on the projected low-energy bands, with the remote bands entering only through $r_0$ as a static background. The RPA is solved in one and two dimensions with an arbitrary kernel, the inter-unit lattice sums and the transverse form factor, on $q_\parallel$ from $10^{-4}$ to \SI{10}{keV} (200 logarithmic points) and $\omega$ up to $v_D/r_{\mathrm{t}}$ for tubes or \SI{2}{eV} for sheets with at least 2000 points; stacks of period $a_s$ keep $q_z$ unreduced, with at least 64 points per $2\pi/a_s$. The interlayer spacing and the substrate $\varepsilon_b(\omega)$ with its polar phonons complete the input.

A semiconducting tube screens its electron--hole pairs weakly, so its lowest excitations are excitons, which the RPA does not contain. For such a tube the RPA step is followed by a ladder (Bethe--Salpeter) step on the same projected bands, with the interaction of the tube statically screened by the projected response, which gives the exciton energies and their form factors along $q_\parallel$. The lowest exciton energy must be stable to \SI{10}{meV} when the bands and the $q_\parallel$ grid are doubled and must reproduce the measured optical transition of the same tube, and the ladder adds $10^2$ to $10^3$ core-hours per chirality.

The band-projected $f$-sum of the linear-band model, $\int\omega S\,\dd\omega = N_{\mathrm{ch}}v_Dq^2$ per unit length of tube in that normalization, holds to $10^{-3}$. The undoped Dirac sheet reproduces its closed-form bubble to $10^{-4}$, a finite stack of about 50 layers converges to the infinite-stack form to \SI{1}{\percent}, and the Yambo slab calculation with a Coulomb cutoff cross-checks one sheet. A metallic nanotube couples to dark matter through its charge mode only, with all four spin and valley channels; a one-channel model is low by a factor $N_{\mathrm{ch}} = 4$ at fixed Luttinger parameter. For dots, the envelope-function solver (spherical-well levels up to $\ell = 20$ and $n = 10$) and the dielectric-confinement factors are checked against tight-binding or DFT clusters of up to 2000 atoms. The model module runs in less than one core-hour; the DFT bubbles of real sheets cost $10^2$ to $10^3$ core-hours and the dot clusters $10^4$.

\subsection{Ejection factors of graphene and nanotubes}

A plane-wave final state is not orthogonal to the occupied orbital it leaves, and the overlap adds a spurious monopole term to the ejection amplitude. Orthogonalizing the plane wave to the occupied states removes it. For a carbon $2p_z$ orbital of Slater form with the exponent $1.568/a_0$ of Clementi and Raimondi, the ratio of the orthogonalized to the plain squared amplitude ranges from 0.4 to 4.3 at kinetic energies of 50 to \SI{100}{eV} and reaches 740 at \SI{10}{eV} (\cref{sec:16.3}). The large values occur where $q$ far exceeds the final momentum and the plain amplitude has already fallen as $|\vect{p} - \vect{q}|^{-5}$, so a table of such ratios is not a rate correction. A rate needs the factor weighted by the halo,
\begin{equation}
  R_{\mathrm{OPW}}(E_{\mathrm{kin}}) = \frac{\int\dd q\,q\,\eta\,F_{\mathrm{DM}}^2\sum_{\mathrm{occ}}|M_{\mathrm{OPW}}|^2}{\int\dd q\,q\,\eta\,F_{\mathrm{DM}}^2\sum_{\mathrm{occ}}|M_{\mathrm{PW}}|^2} ,
  \label{eq:E.21-ROPW}
\end{equation}
dimensionless, per occupied orbital and in total, for graphene and nanotubes, with $\eta$ the mean inverse speed of the standard halo. The factor is computed at $\mathcal{X}$ = IP with frozen orbitals and a single continuum orthogonalized to all occupied states, the carbon 1s and the $\sigma$ and $\pi$ bands.

The occupied states come from a DFT ground state (SG15) as maximally localized Wannier functions, with the 1s orbital from an all-electron atomic calculation, since the pseudopotential has none. Their plane-wave transforms and mutual overlaps come by FFT, and the integrals of \cref{eq:E.21-ROPW} use $2\times10^4$ final directions on a Fibonacci set and a logarithmic $q$ grid from 1 to \SI{60}{keV}, which covers heavy-mediator kinematics to $m_\chi \approx \SI{10}{MeV}$. Orthogonality holds to $10^{-6}$, and $R_{\mathrm{OPW}}$ must change by less than \SI{1}{\percent} when the directions or the $q$ grid are doubled and the Wannier spread is reminimized; the whole calculation takes CPU minutes. The final-state interaction with the hole is a separate factor, bounded above by the Fermi factor of an unscreened point charge, \cref{eq:2.2-fermi} (7.3, 3.4 and 2.6 at 10, 50 and \SI{100}{eV}), and computed in CPU hours with scattering states of the sheet's Hartree potential plus the screened hole. An orthogonalized plane wave is not a scattering state of the hole, so multiplying it by a Fermi factor double counts. Both factors replace the bounds of \cref{sec:16.3}.

\subsection{Ionic color centers}

In LiF, NaF, MgO and LiH an electronic deposit can be stored promptly as a Frenkel pair: a self-trapped exciton converts within picoseconds into an F center and an H center. The yield factorizes as in \cref{sec:16.6}, and two of its factors have not been computed: the probability $p_{\mathrm{FH}}(T)$ that a self-trapped exciton becomes a pair at millikelvin, and the survival $p_{\mathrm{surv}}$ of the close pair against tunneling recombination. The calculation gives these two probabilities, the direction-resolved displacement threshold $E_d(\hat{\vect{r}})$ of the anion sublattice, the charge levels that set $p_q$ and the F-band cross-sections, tabulated as yields per recoil energy. On $\mathfrak{T}$ the nuclear-recoil seed is in the impulse regime of $\mathcal{K}$ and the electronic seed at $\mathcal{X}$ = RPA with local fields; the yield factors themselves lie outside $\mathfrak{T}$, as properties of the final state, and come from the dynamics below.

The adiabatic surfaces of the self-trapped exciton and the pair come from constrained-occupation hybrid DFT ($\Delta$SCF with HSE or PBE0, the hybrid built on the functional of Perdew, Burke and Ernzerhof; exact-exchange fraction 0.25 and 0.40, both reported) in 64- to 216-atom cells, large enough for pair separations up to the third neighbor (\cref{app:E.10}). Direct surface hopping on hybrid surfaces would cost $10^6$ to $10^7$ core-hours, so a machine-learned surface is fitted to the $\Delta$SCF points and validated on held-out ones to \SI{50}{meV}, and decoherence-corrected surface hopping runs on it with at least $10^3$ trajectories, which fixes $p_{\mathrm{FH}}$ to \SI{10}{\percent}. The self-trapped-exciton and F--H pair energies are converged to \SI{0.1}{eV} between the 64- and 216-atom cells. The threshold $E_d(\hat{\vect{r}})$ comes from recoil molecular dynamics (\cref{app:E.11}) over $10^4$ zero-point-sampled directions, converged to \SI{1}{eV} when the directions are doubled, and $p_{\mathrm{surv}}$ from the tunneling-path calculation of \cref{app:E.22} applied to the close pair. The F-band cross-section comes from a BSE on the defect supercell. The charged components need the finite-size corrections of \cref{app:E.10}, and at millikelvin only the athermal channel acts. The $\Delta$SCF points dominate a cost of about $10^5$ core-hours per material. The calculation has a direct control: tagged neutron scattering on LiF at \SI{10}{mK} followed by single-center optical readout measures the product of the recoil factors, and \SI{5.9}{keV} X-rays measure $p_{\mathrm{FH}}$ through the quenching of the exciton luminescence.

\subsection{Bond breaking in molecular solids}

A nuclear recoil on one atom of a molecule splits, for an internal coordinate independent of the intermolecular potential, into a center-of-mass recoil of the molecule and an internal excitation at the reduced momentum $(\mu_{12}/M_1)\vect{q}$, with $\mu_{12}$ the reduced mass of the bond and $M_1$ the mass of the struck atom. The response is their convolution (\cref{sec:16.6}), dissociation is the internal response above the bond energy, and in the crystal the fragments also have to escape the cage. The first target is anthracene, the molecular crystal above, with a C--H bond as the internal coordinate, and the calculation gives the internal response $S_{\mathrm{int}}$, the center-of-mass response $S_{\mathrm{CM}}$ and the escape probability $p_{\mathrm{esc}}$ per internal energy and direction. The internal motion is in the impulse regime of $\mathcal{K}$ with a bound center of mass; at the relevant momenta, about \SI{e2}{keV}, the recoil-free fraction is essentially zero, so the center-of-mass recoil is impulsive with a Doppler width from its zero-point momentum.

The internal response comes from Morse or ab initio potential curves, the center-of-mass response from phonopy or path-integral moments (\cref{app:E.8}), and $p_{\mathrm{esc}}$ from molecular dynamics with a machine-learned potential whose training set contains the dissociative region (\cref{app:E.11}). With $10^3$ trajectories per internal energy and direction, $p_{\mathrm{esc}}$ is known to about \SI{3}{\percent}, for about $10^3$ core-hours in all. In the free-molecule limit, $S_{\mathrm{CM}}$ replaced by free recoil and $p_{\mathrm{esc}} = 1$, the calculation must reproduce the dissociation rate of Essig et al.~\cite{essig2017detection}, and $S_{\mathrm{CM}}$ must satisfy the $f$-sum for the whole molecular mass. The photodissociation quantum yield of the same molecular crystal at millikelvin, with fragments detected by infrared absorption, measures $p_{\mathrm{esc}}$.

\begin{table}[t]
  \centering
  \caption{Dopant, color-center, molecular and low-dimensional calculations: targets, main parameters, acceptance tests and cost in core-hours. CNT, carbon nanotube; ML, machine-learned.}
  \label{tab:E21-params}
  \small
  \setlength{\tabcolsep}{4pt}
  \begin{tabular}{p{0.18\textwidth}p{0.15\textwidth}p{0.23\textwidth}p{0.2\textwidth}p{0.09\textwidth}}
    \toprule
    quantity & targets & parameters & acceptance & cost \\
    \midrule
    $\langle S^D(\vect{q},\omega)\rangle$, $\mathcal{U}_{j'j}(\vect{G})$ & Si:P & $10^4$--$10^5$ sites; 100 realizations & sum per donor; dilute limit & $10^3$--$10^4$ \\
    $\Delta S$, $\theta^2_{\mathrm{CT}}$, tables & anthracene, t-stilbene & $4^3$, 200 bands; 590 directions & partial $f$-sum \SI{2}{\percent} & $10^3$--$10^4$ \\
    $S$ of tubes, sheets, stacks, dots; $\tilde R^{\mathrm{2D}}$, $r_0$ & graphene, CNT, stacks, dots & 200 $q_\parallel$; ${\ge}2000$ $\omega$ & projected $f$-sum $10^{-3}$ & ${<}1$ to $10^4$ \\
    tube excitons & semiconducting CNT & ladder on projected bands & lowest exciton \SI{10}{meV} & $10^2$--$10^3$ \\
    $R_{\mathrm{OPW}}$ & graphene, CNT & $2\times10^4$ directions; $q$ 1--\SI{60}{keV} & orthogonality $10^{-6}$ & minutes \\
    $p_{\mathrm{FH}}$, $p_{\mathrm{surv}}$, $E_d(\hat{\vect{r}})$ & LiF, NaF, MgO, LiH & 64--216 atoms; ML surface; ${\ge}10^3$ trajectories & tagged neutrons, X-ray quenching & $10^5$ \\
    $S_{\mathrm{int}}$, $S_{\mathrm{CM}}$, $p_{\mathrm{esc}}$ & anthracene & $10^3$ trajectories per point & free-molecule limit & $10^3$ \\
    \bottomrule
  \end{tabular}
\end{table}
